\documentclass[a4paper,12pt]{article}
\usepackage[utf8]{inputenc}
\usepackage[spanish]{babel}
\usepackage{amsmath}
\usepackage{amsfonts}
\usepackage{amssymb}
\usepackage{siunitx}
\usepackage{geometry}
\usepackage{enumitem}

\geometry{margin=2.5cm}

\title{Examen de Física I - Forma A}
\author{Nombre: \underline{\hspace{8cm}} \quad Matrícula: \underline{\hspace{3cm}}}
\date{Fecha: \underline{\hspace{4cm}}}

\begin{document}

\maketitle

\section{Parte I: Teoría (4 puntos - 1 pt c/u)}
Responda brevemente a las siguientes preguntas:

\begin{enumerate}[itemsep=2em]
    \item \textbf{Interdisciplinariedad}: Mencione tres ciencias con las que la física tiene una relación directa.
    \item \textbf{Cifras Significativas}: Explique la regla de los ceros situados a la izquierda y la de los ceros intermedios.
    \item \textbf{Divisiones}: Mencione las dos grandes divisiones de la física y dé tres ejemplos de ramas de la Física Clásica.
    \item \textbf{Unidades SI}: Enumere cinco (5) magnitudes fundamentales del Sistema Internacional junto con sus unidades básicas.
\end{enumerate}

\vspace{2cm}

\section{Parte II: Práctica (6 puntos - 1.5 pts c/u)}
Resuelva los siguientes problemas mostrando todo su procedimiento. Redondee según las reglas de cifras significativas aprendidas.

\begin{enumerate}[itemsep=4em]
    \item \textbf{Conversión de Longitud}: Convierta \SI{9.5}{\km} a pies (ft). (Use $1\text{ m} = 3.28084\text{ ft}$).
    \item \textbf{Conversión de Velocidad}: Convierta \SI{12}{\mile\per\hour} a \si[per-mode=symbol]{\metre\per\second}.
    \item \textbf{Notación Científica}: Realice la siguiente operación y exprese el resultado en notación científica: $(4.5 \times 10^7) \times (1.8 \times 10^{10})$.
    \item \textbf{Cifras Significativas}: Realice la siguiente suma y redondee correctamente: $15.482 + 3.1 + 102.55$.
\end{enumerate}

\end{document}
